\section{Differentiable coordinates and transverse directions}
\label{sec:differentiability}

We now convert the crossing measures into local geometric information.
The input is the countable family $\mathcal D$ of dyadic time windows,
the edge endpoint marginal $\lambda$ on the rank cube $Z$, its conditional time
probabilities $h_z$, and the crossing measures $\sigma_{CD}$ constructed
above.  We use three established facts: $h_z$ is nonatomic for
$\lambda$-almost every $z$;
\begin{equation}
 \lambda_J=h_z(J)\lambda,
 \qquad
 \lambda_J\ll\sum_{\substack{C,D\subset J\\ C\cap D=\varnothing}}
                 \sigma_{CD};
 \label{eq:analytic-input-ac}
\end{equation}
and the $(C,D)$-coordinate projection of $\sigma_{CD}$ is dominated by
planar Lebesgue measure.  The sum in \eqref{eq:analytic-input-ac} specifies
null sets; it may equivalently be replaced by a finite measure formed
from a countable sum with strictly positive weights.  Our objective is
to find a positive part of $\lambda$ on which
three different sweep directions have three different tangent directions.

\subsection{Calculus for the limiting coordinates}

We first record the two analytic facts needed below. Differentiability
of coordinatewise monotone functions is classical
\cite{ChabrillacCrouzeix1987}; the precise finite-dimensional statement
is also given in \cite[Theorem~1(f)]{BorweinBurkeLewis2004}.
We include the argument needed here. Monotonicity in
each coordinate may have either prescribed sign; reflecting a coordinate
reduces the first lemma to the nondecreasing case.

\begin{lemma}\label{lem:monotone-differentiability}
A bounded Borel function on an open rectangle in $\mathbb R^2$ that is
monotone in each coordinate is Fr\'echet differentiable almost everywhere
with respect to planar Lebesgue measure.
\end{lemma}

\begin{proof}
It suffices to work in a compact subrectangle and assume that both
monotonicities are nondecreasing.  Slice monotonicity shows that the
distributional derivatives $D_1f,D_2f$ are locally finite positive
measures.  Put $\nu=D_1f+D_2f$ and
$Q_r(x)=x+[-r,r]^2$.  For a square lying sufficiently far inside the
domain, monotonicity and integration over
$x+[-2r,-r]^2$ give
\begin{equation}
 \operatorname{osc}_{Q_r(x)}f
 \le f(x+(r,r))-f(x-(r,r))
 \le \frac{\nu(Q_{3r}(x))}{r}.
 \label{eq:monotone-oscillation}
\end{equation}
Indeed, compare every point $u$ of that lower square with
$u+(3r,3r)$.  The difference of their function values is at least the
middle difference in \eqref{eq:monotone-oscillation}.  Integrating the
horizontal increment and then the vertical increment by slices bounds
the resulting integral by $r\nu(Q_{3r}(x))$; the lower square has area
$r^2$.

The upper area density of $\nu$ is finite at almost every $x$.
Consequently, at almost every $x$ there are $M_x,r_x>0$ such that
$|f(y)-f(x)|\le M_x|y-x|$ whenever $|y-x|<r_x$.
These points are covered by countably many measurable sets $E_m$ on
which the restriction of $f$ is Lipschitz.  For example, on a fixed
bounded subrectangle one may impose a uniform bound on the corner
differences in \eqref{eq:monotone-oscillation} for every rational
$0<r<1/m$; boundedness deals with pairs at distance at least $1/m$.
Extend $f|_{E_m}$ to a Lipschitz function $g_m$ on the plane.

By Rademacher's theorem~\cite{Rademacher1919}, in the form stated in
\cite[Chapter~2, Theorem~1.4]{Simon2014GMT}, the Lipschitz extension
$g_m$ is differentiable almost everywhere. Thus, at almost every point
$x\in E_m$, the set $E_m$ has density one and $g_m$ is differentiable.  Its first-order expansion also holds for
$f$ at arbitrary nearby points.  To see this, write $y=x+v$ and choose
$y^-,y^+\in E_m$ with
\[
 y^-_i\le y_i\le y^+_i\quad(i=1,2),
 \qquad |y^\pm-y|=o(|v|).
\]
Such points exist by density one: a missing box of side comparable to
$|v|$ in either prescribed corner would contradict density, and the
relative box size can tend to zero.  Monotonicity squeezes $f(y)$ between
$g_m(y^-)$ and $g_m(y^+)$, giving
$f(x+v)=f(x)+Dg_m(x)v+o(|v|)$.
Countably many sets $E_m$ cover almost all points, proving the claim.
\end{proof}

The next lemma concerns the entire set of differentiability points with
singular derivative.  This distinction matters because a crossing measure
need not be absolutely continuous in the coordinates of the domain. The
statement is the equal-dimensional critical-image estimate; a matching
formulation appears in \cite[Appendix~B, Corollary~B.4]{PayneRedaelli2025preprint}.
We give the covering proof for the precise hypotheses used here.

\begin{lemma}\label{lem:critical-image-null}
Let $F:\Omega\to\mathbb R^2$ be any map, where $\Omega\subset\mathbb R^2$
is open.  The image of the set of points at which $F$ is Fr\'echet
differentiable and $\operatorname{rank}DF<2$ has planar outer measure zero.
\end{lemma}

\begin{proof}
Restrict that set to $S\subset B(0,R)$ where $\|DF(x)\|\le M$.
Fix $\varepsilon>0$.  For each $x\in S$ choose arbitrarily small radii
$r<1$ such that $B(x,5r)\subset\Omega$ and
\[
 |F(y)-F(x)-DF(x)(y-x)|\le\varepsilon|y-x|
 \quad\text{for }y\in B(x,5r).
\]
The elementary $5r$ covering lemma gives disjoint balls $B(x_i,r_i)$
whose fivefold enlargements cover $S$.  Since their centers lie in
$B(0,R)$, we have $\sum_i r_i^2\le (R+1)^2$.
The set $F(B(x_i,5r_i))$ lies in the $5\varepsilon r_i$ neighborhood
of a segment of length at most $10Mr_i$.  Its outer area is therefore
at most $C(M\varepsilon+\varepsilon^2)r_i^2$.
Summing these covers and letting $\varepsilon\downarrow0$ proves
that $F(S)$ is null.  A countable union over $R,M$ completes the proof.
\end{proof}

\subsection{A positive part with independent gradients}

Because $h_z$ is a nonatomic probability, for almost every $z$ it gives
positive mass to two dyadic windows whose closures are disjoint and
contained in $(0,1)$.  There are only countably many pairs, so we can
fix such windows $K,K'$ for which
\[
 E=\{z:h_z(K)>0,\ h_z(K')>0\}
\]
has positive $\lambda$-measure.  Equation
\eqref{eq:analytic-input-ac} for $K'$ implies that $\lambda|_E$ is
absolutely continuous with respect to the sum of its crossing measures.
We may further restrict $E$ to a positive-measure set on which
\begin{equation}
 \lambda|_E\ll\sigma_{AB}
 \quad\text{for some disjoint }A,B\subset K'.
 \label{eq:one-crossing-part}
\end{equation}
For completeness, form a finite measure $\Sigma$ from the countable sum
by assigning strictly positive weights to its summands.
If $g_{AB}=d\sigma_{AB}/d\Sigma$, the sets $\{g_{AB}>0\}$ cover
$\Sigma$-almost all points.  One of them has positive
$\lambda|_E$-measure; restriction to it proves
\eqref{eq:one-crossing-part}.
Thus the projection of $\lambda|_E$ to
\[
 p=(z_A,z_B)\in(0,1)^2
\]
is absolutely continuous with respect to planar Lebesgue measure.
The boundary of the square has zero mass by the same projection bound.

For each $C\subset K$, express the finite coordinate $r_C$ as a function
of $(r_A,r_B)$.  These functions are bounded and coordinatewise monotone,
with signs determined by the relative order of the windows.  Extract a
subsequence on a countable dense grid, simultaneously for all $C$.
The lower monotone extensions of the grid limits give Borel monotone
functions $f_C$.  At every continuity point $p$ of $f_C$, squeezing
between grid points below and above $p$ shows convergence even along
moving arguments $p_m\to p$.

To apply this convergence to the endpoint measure, take
$z\in\operatorname{supp}\lambda$ with pair coordinate $p\in(0,1)^2$.
There are actual finite endpoint labels $z_m\to z$ along a subsequence:
every open neighborhood of $z$ has positive limiting mass and therefore
contains finite endpoint labels at all sufficiently late indices.
Their pair coordinates tend to $p$, so moving-argument convergence gives,
whenever $f_C$ is continuous at $p$,
\begin{equation}
 z_C=f_C(p).
 \label{eq:coordinate-identities}
\end{equation}
By Lemma~\ref{lem:monotone-differentiability} and the absolute continuity
of the pair projection, we can discard a $\lambda|_E$-null set so that
all these identities hold and all the $f_C$ are differentiable at $p$.

The fixed pair $A,B$ supplies a coordinate chart in which every $f_C$
is differentiable at the pair coordinate of each remaining label.
We now use crossing
measures from pairs $C,D$ inside each positively active window
$J\subset K$ to find independent gradients.  Those crossing measures
are controlled in their own $(C,D)$ coordinates; the critical-image
lemma connects this control to derivatives in the $(A,B)$ chart.

\begin{lemma}\label{lem:independent-gradients}
After discarding a further $\lambda|_E$-null set, every label $z\in E$
has the following property: if $J\subset K$ is dyadic and $h_z(J)>0$,
then some disjoint dyadic $C,D\subset J$ satisfy
\[
 \det\begin{pmatrix}\nabla f_C(p)\\ \nabla f_D(p)\end{pmatrix}\ne0.
\]
\end{lemma}

\begin{proof}
Fix $J$.  Let $B_J\subset E$ be the labels with $h_z(J)>0$ for which
every such determinant vanishes.  For each disjoint $C,D\subset J$,
Lemma~\ref{lem:critical-image-null}, applied to $F=(f_C,f_D)$, gives a
Borel planar null set containing
\[
 \{(f_C(p),f_D(p)):p\text{ is a differentiability point of }F,
                  \ \det DF(p)=0\}.
\]
By \eqref{eq:coordinate-identities}, the $(C,D)$ projection of $B_J$
lies in that null set.  Projection domination therefore gives
$\sigma_{CD}(B_J)=0$.  Equation \eqref{eq:analytic-input-ac} implies
$\lambda_J(B_J)=0$.  Since
$\lambda_J(B_J)=\int_{B_J}h_z(J)\,d\lambda(z)$ and $h_z(J)>0$ on
$B_J$, we get $\lambda(B_J)=0$.  Discard these null sets for all
countably many $J$.
\end{proof}

\subsection{Allowed tangent directions}

Fix a remaining label $z$ and its pair coordinate $p$.
At small rank scales $h$, the preceding convergence and differentiability
have the following iterated consequence: for every fixed $R<\infty$,
\begin{equation}
 r_C=f_C(p)+h\nabla f_C(p)\cdot w+o(h),
 \qquad w=\frac{(r_A,r_B)-p}{h},\quad |w|\le R.
 \label{eq:rank-first-order}
\end{equation}
Here first $h\downarrow0$, and for each fixed $h$ the configuration index
is taken sufficiently large.  The estimate is uniform on the displayed
ball, simultaneously for any fixed finite collection of $C$.
Indeed, differentiate $f_C$ at $p$, then bracket the ball by a finite
grid of continuity points with mesh $o(h)$ and use monotonicity.
Choosing the indices successively gives the estimate for all $C$ along
any sufficiently fast diagonal.  No rate uniform in $p$ is asserted.

Orient each physical sweep direction by increasing original horizontal
coordinate.  For $t\in K^\circ$, forward displacements on a line of
direction $t$ have a common prescribed sign in each of the two rank
coordinates.  Let $Q$ be the resulting closed signed quadrant.
For each $t\in K^\circ$, define $D_p(t)\subset Q\cap S^1$ by the
linear inequalities
\begin{align*}
 \nabla f_C(p)\cdot d&\ge0
   &&\text{if $C\subset K$ lies strictly before $t$},\\
 \nabla f_C(p)\cdot d&\le0
   &&\text{if $C\subset K$ lies strictly after $t$}.
\end{align*}
``Strictly'' here requires a positive time gap.
It is an intersection of closed halfplanes with the quadrant, restricted
to the unit circle, so it is a closed angular interval when nonempty.

Any pair of forward-ordered points on a line of direction $t$, with
bounded zoom coordinates tending to distinct limits $w,w'$, satisfies
\begin{equation}
 \frac{w'-w}{|w'-w|}\in D_p(t).
 \label{eq:all-pairs-direction}
\end{equation}
This follows from monotonicity of $r_C$ along the line and
\eqref{eq:rank-first-order}.  The base point and the line may vary with
the configuration.  Moreover, $D_p(t)$ is nonempty.  In each configuration
follow the forward line through the physical preimage of $p$ until its
rank image first exits the zoom unit ball.  This exit exists: the rank
map is a homeomorphism onto the open square, and the preimage of a compact
subset is compact.  Along a sufficiently fast diagonal the exit points
have a convergent subsequence, and their unit limiting displacement
belongs to $D_p(t)$.

\begin{lemma}\label{lem:separated-tangent-directions}
If $s<t$ lie in $K^\circ$ and $h_z((s,t))>0$, then
$D_p(s)\cap D_p(t)=\varnothing$.  There are three times
$\tau_1<\tau_2<\tau_3$ in $K^\circ$ for which these intervals are
singletons with distinct projective directions.
\end{lemma}

\begin{proof}
Choose a dyadic $J$ whose closure lies in $(s,t)$ and for which
$h_z(J)>0$.  Lemma~\ref{lem:independent-gradients} provides disjoint
$C,D\subset J$ with independent gradients.  A vector belonging to both
$D_p(s)$ and $D_p(t)$ would have zero scalar product with both gradients,
and hence be zero, contrary to unit length.

Consider the support of the nonatomic measure $h_z|_{K^\circ}$.
Remove the endpoints of its complementary intervals and the endpoints
of $K$.  Only countably many points are removed, and any two distinct
remaining support points have positive measure strictly between them.
Their intervals $D_p(t)$ are consequently pairwise disjoint.  At most
countably many disjoint angular intervals can have positive length, since
each contains a different rational angle.  Removing those times still
leaves full $h_z|_K$-measure.  Three remaining times give the claim.
All directions lie in the same quadrant, so distinct unit directions
are also projectively distinct.
\end{proof}

The singleton times in this lemma may depend on $z$.  We also need one
fixed triple for a common normalization of the configurations.  For
every remaining $z$, nonatomicity and $h_z(K)>0$ allow rational times
$t_1<t_2<t_3$ in $K^\circ$ such that both
$h_z((t_1,t_2))$ and $h_z((t_2,t_3))$ are positive.  Countable selection
fixes one triple on a positive $\lambda$-measure subset.  A further
countable selection fixes a dyadic $K_0$ with closure in $(t_1,t_2)$ and
$h_z(K_0)>0$ throughout a positive-measure subset, again denoted $E$.
Thus
\[
 \lambda_{K_0}(E)>0,
 \qquad
 D_p(t_1),\ D_p(t_2),\ D_p(t_3)
 \text{ are pairwise disjoint compact angular intervals.}
\]
All restrictions made here are restrictions of the original spatial
measure $\lambda$.  Positive conditional time mass was used only in
the stated measurable events; positive spatial mass follows from the
countable selections.  These conclusions supply the directional input
for the local geometric argument.
